% !TEX root = ../main.tex

\section{Concluding Summary}\label{sec:CS_ch2}

In this chapter, the development of the FW-H method has been described. Test cases have been developed to validate sources in stationary and motion. Both the thickness and loading noise, which are the sources that are evaluated in the permeable surface formulation, were validated whilst testing the limits of time integration algorithms, discretisation resolutions, and analytical assumptions. Overall, the FW-H code has been shown to be able to predict these analytical sound sources with a sufficient accuracy.

The monopole validation have shown the numerical FW-H code's range of capability in predicting the thickness sources, using cases such as co-stationary, co-translation, and a translating source (Figures \ref{fig:monopole_SS_RT}-\ref{fig:monopole_TS_AT}). The code was able to predict phenomena such as phase lagging, convective amplification, and Doppler shifting. The validation varied the discretisation resolutions between 275, 100, and 35 grid points per wavelength and 30, 15, and 5 time steps per period; with which the numerical FW-H achieved a decreasing errors between $\sim0.5\%$ and $\sim30\%$ using the finest and coarsest resolutions, respectively. This shows that the finest resolution is sufficient to properly predict these simple cases.

Alongside the monopole validation, the Retarded Time and Advanced Time algorithms have been successfully validated. Given a sufficient spatial and temporal discretisation resolutions, both algorithms achieved good agreements with the analytical solutions across all test cases. While the Retarded Time is more straightforward to implement and validate, the Advanced Time algorithm is more advantageous when applied to real CFD cases as it allows the use of the source's non-uniform time-stepping while being compatible to parallel computing \cite{Brentner1997, casalino-2001}.

% % MONOOPOLE DISCUSSION %
% In the grid resolution comparisons, while the finest resolution of $40\times40$ has very good agreement with the analytical solutions at less than $1\%$ error in all cases, the errors progressively increase to $\sim4\%$ and $\sim 30\%$, for coarser resolutions. The plots show at lower resolution, the numerical results could not capture the peaks and troughs of the harmonic wave, causing dissipation. Given the acoustic wavelength is $\sim214m$ (from the prescribed frequency and speed of sound values), the coarsest grid has an arc distance between two nodes of 1.25 meters. The errors shown in the grid sensitivities are mainly due to geometric error of the spherical domain, rather than the grids not being able to capture the wavelengths. This shows that this grid resolution is sufficient given the input and control surface radius.

% The time step sensitivity study exhibits similar trends. The $dt=0.005s$ and $dt=0.025s$ sizes yield errors at less than $1\%$ and $2\%$, respectively. However, largest time-step size of $dt=0.125s$ (an $8Hz$ sampling frequency), is insufficient to capture the input acoustic frequency of $10rad/s$ or $\sim1.6Hz$, as it can only sample at five times the input. This can be seen in the lags in the plots as the numerical solution attempted to keep up with the wave. It should be noted, however, that the spatial discretisation error that come from these setups mainly stems from geometric approximation error, rather than the grids not being able to resolve the wavelengths.

Moving to the dipole validation, while the direct analytical dipole derivations were limited, the numerical FW-H predicted the loading sources in good agreement with the respective assumption's analytical solutions, provided the test cases follow the assumption. The stationary assumption fails when the source is in motion, which is due to the assumption made in its derivation where $\partial r/\partial\tau=0$, which causes the numerical FW-H to introduce a phase lag in its prediction (Figure \ref{fig:dipole_validation_stationary}). As shown in Equation \ref{eq:dipole_u1}, there is no Doppler factor $1-M_r$ in any of the terms, which causes any motion to be invalid. On the other hand, the plane wave assumption is invalid when the permeable surface is placed in the near-field, which causes the numerical solution to over-predict the amplitudes while introducing a phase lag (Figure \ref{fig:dipole_validation_planewave}). This is because the velocity field is assumed to be in-phase with the pressure field, which is not the case in the acoustic near-field.

In contrast to the analytical derivations, the dual-monopole method is significantly more robust in all cases. By superimposing two equal monopoles of opposite signs separated at a distance of $k\delta\ll1$ ($k$ is the wave number), the code instead uses the analytical solution of the monopole, thereby omitting the assumptions made in the dipole velocity derivations. Even so, although the analytical dipole solutions were limited, they provided a good baseline against which the dual monopole approach was used to validate the loading noise. Additionally, the dual monopole approach allows the visualisation of the sound directivity, which is highly relevant in the following chapters to show how the sound waves are distributed in space.

To conclude, the FW-H method has been rigorously validated across different cases while simultaneously testing the discretisation resolution and time integration method. Validation strategies presented in this chapter are the first steps in building confidence in the FW-H code's noise prediction capability. The next chapter verifies the FW-H code's accuracy when used to mimic sound signatures coming from a known source as well as its sensitivity to varying Mach numbers, which is to understand how discretisation parameters affect its prediction when used in more complex cases.

% DIPOLE DIDSCUSSION %


% For the zero relative Mach number assumption, the tests examine how relative motion affect the results. In each case, the observer is put at $[50,50,0]$. To examine the integration accuracy, the first case examines a stationary problem, and the second one a relative motion problem, moving the source at $150m/s$ along the x-direction. In the plane wave assumption, the tests only examine the relative motion and putting the observer at $[500,500,0]$, to be in the "far-field." The control surface radius is expanded to $r=400m$ in the valid case and to $r=1m$ in the invalid case.

% For stationary source and true far-field observer (Figures \ref{fig:near_valid} and \ref{fig:far_valid}), the results have good agreements at $<6\%$ error. However, introducing translation into the former and setting near-field surface to the latter violate the assumptions. For the $M_r=0$ case, the errors are due to the non-existent relative motion in Equation \ref{eq:dipole_u1}; for the plane wave assumption, the control surface being in the near-field of the source. The agreement as the source moves away from the observer is due to the observer eventually being in the physical far-field.

% In contrast to the analytical derivations, the dual-monopole method is significantly more robust in all cases, with highest error at $<8\%$. By superimposing two monopoles separated by $k\delta\ll1$, the solver retains the Doppler factor, which is omitted by the zero relative motion assumption. Similarly for the bottom figures, as the monopole was validated using a control surface in the near-field, the results generally agree in the cases which the far-field assumption failed when using the same control surface. It is noted that in Figure \ref{fig:dual_far_invalid}, the result exhibits more significant errors as the source moves away. This could be attributed to the limitations in the solver and will be discussed in the next chapter.

% Comparing the dipole derivations to the dual monopole approach shows that there are apparent deviations in results other than the stationary case. For Figure \ref{fig:dual_dipole_near_invalid}, it is obvious that the cause is the lack of Doppler effect in the assumption. Figure \ref{fig:dual_dipole_far_valid} is interesting due to a significant difference in amplitude as the sources approach the observers. Then as they move away, a good agreement in amplitude is observed. This is likely due to the imposing of the $x$-axis in the dipole source, because as it moves away, the $y$-axis contribution in the emission distance becomes insignificant.